
\chapter*{ABSTRACT}

\footnotesize{
\begin{description}
  \item[TITLE:] \MakeUppercase{Detection of Diabetic Retinopathy in Fundus Images Using Deep Learning Algorithms.}
  \astfootnote{Bachelor Thesis}
  \item[AUTHOR:] {\autora}, {\autorb}
  \asttfootnote{\facultad. \escuela. Advisor: \director. Co-advisor: \codirectora}
  \item[KEYWORDS:] DIABETIC RETINOPATHY, CONVOLUTIONAL NEURAL NETWORKS, DEEP LEARNING, FUNDUS IMAGES, DIGITAL IMAGE PROCESSING, TRANSFER LEARNING.
  \item[DESCRIPTION:]\hfill \\ Diabetic retinopathy (DR) is a microvascular complication of diabetes that damages the blood vessels in the retina, affecting 145 million people worldwide and constituting one of the main causes of preventable blindness. In Colombia, with a ratio of 1.56 ophthalmologists per 100,000 inhabitants and approximately 1.2 million people with diabetic retinopathy, automated detection represents a critical solution to expand screening coverage. Conventional diagnosis through fundus examinations is costly and limited by the availability of specialists, especially in regions with scarce medical resources. This work develops models based on convolutional neural networks (CNNs) to detect the disease and classify its severity levels through binary classifiers (presence/absence) and multiclass classifiers (five levels). Transfer learning was applied using MobileNetV2, ResNet50V2, InceptionV3, and VGG19 architectures, trained with 35,121 images from the Kaggle public dataset. Preprocessing included color normalization, peripheral masking, class balancing through subsampling and data augmentation, incorporating Multi-Head Attention blocks, GaussianNoise layers, and regularization techniques (Dropout, L2, early stopping). The model was evaluated using sensitivity, precision, F1-score, and AUC-ROC, with MobileNetV2 with preprocessing through filtering and optimized threshold (0.4826) being the best-performing model, achieving 96.5\% sensitivity in binary classification with bootstrap confidence intervals (95\% CI: AUC-ROC 0.9636–0.9936). For multiclass classification, 55\% accuracy and 65\% sensitivity were achieved, identifying the green channel as most effective due to its better contrast for hemorrhagic lesions. External validation on MESSIDOR-2, IDRiD, and FOSCAL datasets confirmed the model's generalization capability, demonstrating its viability as a support tool for clinical screening of diabetic retinopathy.

\end{description}}\normalsize
% ------------------------------------------------------------------------
